% Part: first-order-logic
% Chapter: syntax
% Section: intro-syntax

\documentclass[../../../include/open-logic-section]{subfiles}

\begin{document}

\olfileid{fol}{syn}{itx}

\olsection{پېژندنه}

د لومړۍ درجې منطق د تيورۍ او ميتاتيورۍ د پراختيا دپاره بايد تر هر څه
مخکښې د دې منطق د عبارتونو نحوه او معنٰی پوهنه تعريف کړو۔ د لومړۍ درجې
منطق عبارتونه ترمونه او !!{formula}s دي۔ ترمونه له !!{variable}s،
!!{constant}s او !!{function}s څخه جوړېږي۔ !!^{formula}s بيا له
!!{predicate}s او ترمونو څخه جوړېږي (په دې ډول تر ټولو واړۀ، يعنې
``اټومي'' !!{formula}s جوړېږي)، او بيا له اټومي !!{formula}s څخه د
منطقي تړونکو او کميت ټاکونکو په وسيله نور پېچلي فارمولونه جوړولے شو۔ د
جوړېدو د قاعدو د ټاکلو ډېرې بېلابېلې لارې شته؛ موږ يوازې يوه ممکنه لاره
وړاندې کوو۔ نور نظامونه بېل سمبولونه غوره کوي، د تړونکو بېلې ټولګې د
لومړنيو په توګه ټاکي، او قوسونه په بله طريقه کاروي (يا، لکه د تش په نامه
پولنډۍ ليکنې په حالت کښې، ئې بيخي نۀ کاروي)۔ خو ټولې لارې په دې کښې
شريکې دي چې د جوړېدو قواعد د ترمونو او !!{formula}s ټولګه
\emph{په استقرايي ډول} تعريفوي۔ که دا کار سم وشي، نو هر عبارت د جوړېدو
د قواعدو له مخې په بنسټيز ډول يوازې په يوه طريقه تر لاسه کېږي۔ هغه
استقرايي تعريف چې \emph{يوازېنۍ لوستنې وړ} عبارتونه رامنځته کوي، موږ ته
دا امکان راکوي چې همدې طريقې---يعنې استقرايي تعريف---سره دغو عبارتونو ته
معناوې هم ورکړو۔

\end{document}
